% !PRINTED 629
\weberpage{629}
%<<CONT
\[
\begin{aligned}
\pi_1 \pi_2 &= (1, 2, 6) (3, 4, 5) \\
\pi_1 \pi_2 \pi_3 &= (1, 6, 5, 4) \\
\pi_1 \pi_2 \pi_3 \pi_4 &= (2, 4, 6, 3, 5) \\
(\pi_1 \pi_2 \pi_3 \pi_4)^3 &= (2, 3, 4, 5, 6).
\end{aligned}
\]
Hiernach lässt sich leicht eine zur Gruppe $C$ gehörige Function der sechs Grössen $v$ bilden.

Eine Function, wie
\[
v_1 v_3 + v_2 v_4 + v_5 v_6,
\]
bleibt durch $\pi_1$ und durch $\pi_1 \pi_2$ ungeändert; wenden wir darauf aber die Potenzen der cyklischen Permutation $(2, 3, 4, 5, 6)$ an, so gehen daraus fünf Functionen hervor, die wir so bezeichnen wollen:
\weq{2}{
\begin{aligned}
% [sic] Print has a prime after $v_4$ in this row; the parallel combinatorial pattern requires unprimed $v_4$.
w_0 &= v_1 v_2 + v_4' v_5 + v_3 v_6 \\
w_1 &= v_1 v_4 + v_2 v_6 + v_3 v_5 \\
w_2 &= v_1 v_6 + v_2 v_5 + v_3 v_4 \\
w_3 &= v_1 v_3 + v_2 v_4 + v_5 v_6 \\
w_4 &= v_1 v_5 + v_2 v_3 + v_4 v_6.
\end{aligned}
}

Wenn wir auf die Functionen (2) die Permutationen $\pi_1, \pi_2, \pi_3, \pi_4$ anwenden, so gehen diese Functionen nur in einander über, und zwar in folgender Weise:
\[
\pi_1 = (w_0, w_1), \quad \pi_2 = (w_0, w_2), \quad \pi_3 = (w_0, w_3), \quad \pi_4 = (w_0, w_4),
\]
d. h. die $\pi_1, \pi_2, \pi_3, \pi_4$ entsprechen den Transpositionen $(0, 1)$, $(0, 2)$, $(0, 3)$, $(0, 4)$ unter den Indices der $w$, und also entspricht nach \S.\,153, 2. der ganzen Gruppe $C$ die symmetrische Gruppe der Indices der $w$. Eine symmetrische Function der fünf Grössen $w$ die nicht zugleich eine symmetrische Function der $v$ ist, wie die Summe der $w$, ist also eine zu $C$ gehörige Function.

Man kann z. B. die Summe der Quadrate dafür nehmen
\[
W = w_0^2 + w_1^2 + w_2^2 + w_3^2 + w_4^2,
\]
oder eine Function
\[
(\lambda - w_0) (\lambda - w_1) (\lambda - w_2) (\lambda - w_3) (\lambda - w_4)
\]
für ein beliebiges rationales $\lambda$.

Eine solche Grösse $W$ ist die Wurzel einer irreduciblen Gleichung $6^{\mathrm{ten}}$ Grades, deren Coëfficienten symmetrische Functionen der sechs Grössen $v$ sind.

\begin{center}
\rule{3cm}{0.4pt}
\end{center}
